%% Example CV - Brucce (Docente Tecsup & Software/AI Engineer)

\documentclass[11pt,a4paper,sans]{moderncv}
\moderncvstyle{banking}
\moderncvcolor{blue}

\usepackage[utf8]{inputenc}
\usepackage{lmodern}
\usepackage[scale=0.80]{geometry}

\microtypesetup{activate=false}

\renewcommand*{\namefont}{\fontsize{26}{28}\bfseries\upshape}
\colorlet{firstnamecolor}{color1}
\colorlet{lastnamecolor}{color1}
\colorlet{namecolor}{color1}
\renewcommand*{\sectionstyle}[1]{{\sectionfont\color{color1}#1}}

\AtEndPreamble{\hypersetup{
    colorlinks=true,
    linkcolor=blue,
    urlcolor=blue,
    pdftitle={Brucce - CV},
    pdfpagemode=UseNone,
}}

% Personal data
\name{Brucce}{Vargas}
\address{Lima / Arequipa}{Perú}{}
\phone[mobile]{+51 987 654 321}
\email{brucce@tecsup.edu.pe}
\extrainfo{\href{https://github.com/BrucceVT}{GitHub: BrucceVT}}

\begin{document}

\makecvtitle
\vspace{-12pt}

% ============================================================
%     PROFILE STATEMENT
% ============================================================

\small{Docente de Ingeniería de Software \& AI Architect en Tecsup con amplia experiencia en desarrollo Web Full Stack, arquitecturas agénticas con Python, TypeScript y despliegues de microservicios con Docker y REST APIs. Apasionado por la formación técnica avanzada y la automatización inteligente en producción.}

\vspace{6pt}

% ============================================================
%     CORE COMPETENCIES
% ============================================================

\section{Habilidades Principales}
\begin{itemize}
\item \textbf{Lenguajes de Programación}: Python, TypeScript, JavaScript, SQL, Bash.
\item \textbf{Frameworks \& Arquitectura}: FastAPI, Node.js, React, Agentes de IA, Microservicios, REST APIs.
\item \textbf{Infraestructura \& Herramientas}: Docker, Git, Linux, CI/CD, LuaLaTeX, Postman.
\item \textbf{Docencia \& Liderazgo}: Diseño de planes de estudio, mentoría técnica en Tecsup, metodología STAR.
\end{itemize}

\vspace{6pt}

% ============================================================
%     PROFESSIONAL EXPERIENCE
% ============================================================

\section{Experiencia Profesional}

\cventry{2023--Presente}{Docente de Tecnología e Ingeniería}{Tecsup}{Perú}{}{
\begin{itemize}
    \item Instrucción práctica en arquitectura de software, inteligencia artificial y desarrollo Full Stack.
    \item Diseño de módulos de laboratorio interactivos combinando Docker, Python y control de versiones.
    \item Asesoría técnica a estudiantes en el desarrollo de proyectos de fin de carrera y soluciones de producción.
\end{itemize}
}

\vspace{4pt}

\cventry{2022--Presente}{Consultor de Software \& AI Engineer}{Proyectos Tech}{Remoto / Perú}{}{
\begin{itemize}
    \item Diseño e implementación de microservicios con FastAPI y TypeScript.
    \item Integración de modelos LLM y pipelines agénticos con tolerancia a fallos.
    \item Automatización de procesos de integración continua CI/CD y despliegue en contenedores.
\end{itemize}
}

\vspace{6pt}

% ============================================================
%     EDUCATION
% ============================================================

\section{Educación \& Certificaciones}

\cventry{2018--2022}{Ingeniería de Software / Sistemas}{Tecsup / Universidad}{Perú}{}{Especialidad en Arquitectura Agéntica \& Sistemas Distribuidos.}
\cventry{2024}{Especialización en Agentes de IA \& LLMs}{Certificación Internacional}{Online}{}{Capacitación avanzada en orquestación de LLMs y automatización.}

\end{document}
